% !TEX root = ../Projektdokumentation.tex
\subsection{Implementierung der wichtigsten Funktionen}
\label{sec:ImplementierungKernfunktionen}
Die Umsetzung erfolgte schrittweise entlang der fachlichen Teilbereiche.
Zunächst wurden die Datenklassen für Firmen, Benutzer, Mail-Adressen,
IP-Adressen und Audit-Daten angelegt. Danach folgten die Funktionen für
Validierung, Benutzerverwaltung, Rollenprüfung und die Pflegeprozesse. So
konnte jeder Teilbereich einzeln umgesetzt und getestet werden, bevor er in der
Weboberfläche genutzt wurde.

Die fachliche Logik wurde nicht direkt in den Routen oder Templates
implementiert, sondern in Service-Funktionen ausgelagert. Dadurch greifen
Weboberfläche und API auf dieselben Prüfungen zurück. Für die
Benutzerverwaltung wurden die Rollen Root, Admin und User umgesetzt. Root kann
firmenübergreifend arbeiten, Admin verwaltet Benutzer und Fachobjekte innerhalb
der zugeordneten Firma, und User bleibt auf die eigene Firma und die erlaubten
Pflegefunktionen beschränkt. Die Firmenzuordnung wird in der Datenbank beim
Benutzer gespeichert und nach der Anmeldung über die Session dem aktuellen
Benutzerkontext zugeordnet.

Beim Lesen von Mail- und IP-Adressen werden die sichtbaren Datensätze anhand
dieses angemeldeten Benutzers gefiltert. Vor dem Ändern eines Datensatzes wird
zusätzlich geprüft, ob der Benutzer die betroffene Firma bearbeiten darf.
Dadurch wird die Mandantentrennung nicht nur optisch in der Oberfläche, sondern
auch serverseitig durchgesetzt. Der zugehörige Quelltextauszug ist in
Quelltextauszug~\ref{lst:permissions} im \Anhang{app:CodeAuthentifizierung}
dargestellt, weil dort die Filterung sichtbarer Daten und die Prüfung vor
Änderungen direkt zusammengeführt sind.

Ein Schwerpunkt lag auf der Anmeldung und Absicherung der Benutzerkonten. Die
Anwendung nutzt Session-Login, Passwort-Hashing mit Argon2, Login-Audit sowie
Schutzmechanismen gegen wiederholte Fehlanmeldungen. Damit wird verhindert,
dass sensible Verwaltungsfunktionen ohne ausreichenden Schutz erreichbar sind.
Bei der Anmeldung wird die E-Mail-Adresse zuerst vereinheitlicht. Anschließend
prüft die Anwendung, ob für die IP-Adresse oder das Benutzerkonto bereits eine
Sperre aktiv ist. Erst danach wird das Passwort gegen den gespeicherten Hash
geprüft und bei Erfolg eine neue Session angelegt. Fehlgeschlagene und
erfolgreiche Anmeldeversuche werden protokolliert, damit sicherheitsrelevante
Vorgänge nachvollziehbar bleiben. Die Umsetzung ist in
Quelltextauszug~\ref{lst:auth-session} und
Quelltextauszug~\ref{lst:login-security} im
\Anhang{app:CodeAuthentifizierung} referenziert. Diese Stellen sind relevant,
weil sie den Übergang von einer reinen Eingabeprüfung zur geschützten
Benutzersession zeigen.

Für fachliche Änderungen wurde ein Audit-Log umgesetzt. Beim Anlegen, Ändern
oder Löschen von Firmen, Benutzern, Mail-Adressen und IP-Adressen wird ein
Eintrag mit Benutzerbezug, Aktion und Änderungsdaten geschrieben. Sensible
Felder wie Passwörter werden vor der Protokollierung maskiert. Für ausgewählte
Änderungen werden außerdem die Informationen gespeichert, die für ein
kontrolliertes Zurücksetzen benötigt werden. Dadurch kann nicht nur der aktuelle
Stand gepflegt, sondern auch die Entstehung eines Datenstands nachvollzogen
werden. Die Maskierung sensibler Felder und das Schreiben des Audit-Eintrags
sind in Quelltextauszug~\ref{lst:audit-redaction} und
Quelltextauszug~\ref{lst:audit-create} dargestellt. Der dazugehörige Rollback
ist in Quelltextauszug~\ref{lst:audit-rollback} enthalten.

\subsection{Erzielte Projektergebnisse}
\label{sec:ErzielteProjektergebnisse}
Mit Abschluss der Implementierung lag eine lauffähige Webanwendung vor.
Firmen, Benutzer,
Mail-Adressen und IP-Adressen können zentral verwaltet werden. Fachliche
Änderungen werden im Audit-Log dokumentiert und können für definierte Fälle
über einen Rollback wiederhergestellt werden.

Zusammenfassend wurde damit nicht nur eine neue Eingabemaske geschaffen. Der
bisher manuelle Pflegeprozess wurde in eine strukturierte Anwendung mit
Rechtekonzept, Prüfung der Eingaben und Änderungshistorie überführt.

Während der Implementierung zeigte sich, dass die Trennung zwischen
firmenbezogenen und globalen Rechten nicht nur in der Oberfläche
berücksichtigt werden darf. Die Prüfung muss auch in den fachlichen Funktionen
stattfinden, damit geschützte Aktionen nicht über andere Wege ausgeführt werden
können. Deshalb wurde die Rechteprüfung in den Service-Funktionen umgesetzt.
